%---------------------------------------------------------------------
% Front matter: title page, provenance, how to use, exam map.
%---------------------------------------------------------------------
\begin{titlepage}
\thispagestyle{empty}
\begin{tikzpicture}[remember picture, overlay]
  \fill[bmblued] (current page.north west) rectangle
        ([yshift=-9.4cm]current page.north east);
  \fill[bmaccent] ([yshift=-9.4cm]current page.north west) rectangle
        ([yshift=-9.62cm]current page.north east);
\end{tikzpicture}

\vspace*{1.3cm}
{\sffamily\color{white}
  % \rmfamily here: T1/lmss has no small-caps shape, so \scshape inside the
  % sans group would quietly lose the caps.
  {\rmfamily\scshape\lsstyle\large Allama Iqbal Open University, Islamabad}\\[3pt]
  {\footnotesize Department of Mathematics \textbullet\ Faculty of Sciences}\\[24pt]
  {\Huge\bfseries Business Mathematics}\\[10pt]
  {\LARGE\bfseries Solved Question Bank}\\[16pt]
  {\large Course Code \textbf{1429} (also issued as \textbf{5405}) \textbullet\ Units 1--9}\\[6pt]
  {\normalsize BA \textbullet\ B.Com \textbullet\ ADC \textbullet\ BS}
}

\vspace{4.1cm}

\begin{tcolorbox}[enhanced, colback=bmtint, colframe=bmblue, boxrule=0.6pt,
                  arc=3pt, left=14pt, right=14pt, top=12pt, bottom=12pt]
\sffamily\small
\textbf{\color{bmblue}What this document is.} Full worked solutions to the
\saq{self assessment questions} printed at the end of each of the nine units of
the official AIOU course book \emph{Business Mathematics}, Course Code 5405/1429,
Department of Mathematics, AIOU Islamabad (2023 printing).

\medskip
\textbf{\color{bmblue}Why these questions.} AIOU sets its semester assignments and
its final paper from this pool. Solving the book's own end-of-unit questions is
therefore not a proxy for past-paper practice --- it \emph{is} past-paper practice,
one step upstream of the paper.

\medskip
\textbf{\color{bmblue}Status.} Prepared for private study. It is not an AIOU
publication, carries no official endorsement, and must not be copied into an
assignment for submission --- see \hyperref[sec:howtouse]{\emph{How to use this}}
on page~\pageref{sec:howtouse}.
\end{tcolorbox}

\vfill
{\sffamily\footnotesize\color{bmgrey}
Compiled 3 September 2026 \textbullet\ Solutions worked and verified by
back-substitution throughout.\par}
\end{titlepage}

\pagenumbering{roman}
\setcounter{page}{2}

%---------------------------------------------------------------------
\chapter*{Provenance and honest limits}
\addcontentsline{toc}{chapter}{Provenance and honest limits}
\markboth{Provenance and honest limits}{}

\section*{Where the questions came from}

Every question in this document is reproduced from the end-of-unit
\saq{self assessment questions} of the official AIOU course book:

\begin{quote}
\emph{Business Mathematics}, Course Code 5405/1429, Units 1--9.\\
Department of Mathematics, Allama Iqbal Open University, Islamabad.\\
Course Coordinator: Dr Nasir Rehman. Reviewer: Ms Mubashara Hafeez.\\
2023 printing. Distributed by AIOU at
\href{https://online.aiou.edu.pk/LIVE_SITE/SoftBooks/8405.pdf}{online.aiou.edu.pk}.
\end{quote}

I went looking for scanned past papers first. Copies of the Spring 2019,
Spring 2020, Spring 2022 and Autumn 2023 papers for 1429 are circulating on
several study sites, but every one of them is published as a photograph or a
locked PDF with no recoverable text, so none could be transcribed accurately.
The course book, by contrast, is published by AIOU itself with a real text
layer, and its end-of-unit sets are the pool the assignments and the paper are
drawn from. That is the source used here, and it is a better one.

\section*{Two places where the printed book is defective}

The 2023 printing has typesetting faults, and this document does not paper
over them:

\begin{itemize}
  \item \textbf{Unit 6, Question 5} (the electronics materials problem) prints
        identical material requirements for resistors and transistors. That makes
        the coefficient matrix singular and the system inconsistent, so the
        question as printed has no solution. Rather than quietly patch the
        numbers, Section~\ref{sec:u6q5} shows the determinant vanishing, explains
        what a zero determinant means for a production plan, and then solves the
        repaired problem.
  \item \textbf{Unit 8, Question 7} (the production function) prints
        $0.8l^{3}$, which makes $P$ unbounded above, so no maximising $l$ exists.
        Section~\ref{sec:u8q7} works the second-order test on the printed
        function, states plainly that there is no maximum, and then solves the
        intended $0.8l^{2}-0.02l^{3}$ form.
\end{itemize}

A handful of sub-parts in Units 7 and 9 are illegible in the distributed PDF
(radicals and fraction bars collapsed in the text layer). Those are named where
they occur and skipped rather than guessed at.

%---------------------------------------------------------------------
\section*{How to use this}
\label{sec:howtouse}
\addcontentsline{toc}{section}{How to use this}

Read the question. Cover the solution. Work it on paper. \emph{Then} compare.
Reading a worked solution feels like learning and is not; the transfer to an
exam hall comes almost entirely from the failed attempt that precedes it.

Three features are built for that loop:

\begin{description}
  \item[Answer blocks] The final result is boxed in green, so you can check
        whether you got there without reading how.
  \item[Check lines] Most solutions substitute the answer back into the original
        equation. AIOU awards method marks; a substitution check is the cheapest
        way to catch a sign error before the marker does.
  \item[Where marks are lost] The specific, recurring mistake for that question
        type, in red. These are worth more than the solutions.
\end{description}

\begin{pitfall}[On submitting this as an assignment]
Do not. AIOU runs plagiarism checks and these solutions are, by design, in
circulation. The purpose here is to let you check your own working and find the
step where your method breaks. Copy the method, submit your own arithmetic.
\end{pitfall}

%---------------------------------------------------------------------
\section*{What the nine units are worth}
\addcontentsline{toc}{section}{What the nine units are worth}

The paper draws across all nine units, but not evenly. The pattern below is
what the book's own weighting and the circulating past papers agree on.

\begin{center}
\small\sffamily
\begin{tabular}{@{}L{1.1cm} L{5.2cm} L{4.0cm} L{4.4cm}@{}}
\toprule
\textbf{Unit} & \textbf{Topic} & \textbf{Typical exam load} & \textbf{The one thing it tests} \\
\midrule
1 & Probability theory        & Heavy   & Addition vs.\ multiplication rule \\
2 & Random variables          & Medium  & Probabilities must sum to 1 \\
3 & Equations and inequalities& Medium  & Sign flips; two cases for $|\cdot|$ \\
4 & Linear equations          & Heavy   & Slope, then a point \\
5 & Matrices                  & Heavy   & Conformability of the product \\
6 & Determinants and inverses & Heavy   & Cofactor signs; Cramer's rule \\
7 & Differentiation           & Heavy   & Chain rule inside quotient rule \\
8 & Partial derivatives       & Medium  & Hold the other variable constant \\
9 & Optimization              & Heavy   & Second-order test, always \\
\bottomrule
\end{tabular}
\end{center}

\vspace{6pt}
\noindent
Units 4, 5, 6, 7 and 9 between them carry most of the marks and all of the
long questions. If revision time is short, that is the order to spend it in.

\clearpage
\tableofcontents
\clearpage
\pagenumbering{arabic}
\setcounter{page}{1}
